% =====================================================================
% Глава 3. Архитектура программной системы.
% Версия 2 --- после применения правок независимого ревьюера.
% Объем 5--6 страниц.
% =====================================================================

\section{Архитектура программной системы}
\label{sec:architecture}

Программная инфраструктура работы реализована как фреймворк
адаптивных топологий для мульти-агентных систем
(от англ. Adaptive Topologies for Multi-agent systems, ATM)
на языке Python не ниже 3.11, распространяемый как пакет \code{atm}.
Глава описывает архитектуру с акцентом на решения, относящиеся к
научным гипотезам --- адаптивному переключению топологий и
интеграции человека. Исходный код --- в приложении~\ref{app:github}.

Функциональная схема --- на~рисунке~\ref{fig:arch:functional}. Вход ---
описание задачи и YAML-конфиг
(от англ. YAML Ain't Markup Language); выход --- журналы
взаимодействия, метрики и записи о фазах, переключениях топологий
и взаимодействиях с человеком. Между входом и выходом работают
пять блоков --- оркестратор, менеджер фаз и маршрутизаторы, активная
топология с агентами, шлюз HITL и подсистема наблюдаемости.

\begin{figure}[!htbp]
\centering
\fbox{\parbox{0.93\textwidth}{\centering\small
\vspace{0.3cm}
Вход\ \ $\longrightarrow$\ \ Оркестратор эксперимента\ \ $\longrightarrow$\\[0.1cm]
$\big[$\,Менеджер фаз\ \ $+$\ \ Маршрутизатор топологий\,$\big]$\\[0.1cm]
$\longrightarrow$\ \ Активная топология (Star / Chain / Mesh /
Debate / Hierarchical)\\[0.1cm]
$\longleftrightarrow$\ \ Агенты\ \ $\longleftrightarrow$\ \ Инструменты
\& песочница\\[0.1cm]
$\longleftrightarrow$\ \ Шлюз взаимодействия с человеком
(симулятор / интерактивный)\\[0.1cm]
$\longrightarrow$\ \ Подсистема наблюдаемости\ \ $\longrightarrow$\\[0.1cm]
Выход --- метрики, журналы, переходы.
\vspace{0.3cm}
}}
\caption{Функциональная схема предлагаемой мульти-агентной системы.
Прямоугольниками выделены архитектурные блоки фреймворка,
двунаправленные стрелки изображают многократные обмены сообщениями
в пределах одной итерации. Развернутая схема с указанием слоев
агентов, внешних служб и подсистемы хранения приведена в
приложении~\ref{app:uml}.}
\label{fig:arch:functional}
\end{figure}

\subsection{Обзор архитектуры и слои}
\label{sec:arch:overview}

Кодовая база организована в виде тринадцати функционально
изолированных слоев (пакетов высокого уровня), каждый из которых
отвечает за один аспект системы и содержит несколько программных
модулей. Перечень слоев и их назначений приведен в
таблице~\ref{tab:arch:layers}. Такая декомпозиция позволяет
независимо тестировать компоненты, заменять реализации (например,
провайдера языковой модели или шлюз человека в контуре управления
от англ. Human-in-the-Loop, HITL) и добавлять новые топологии без
изменения существующего кода.

\begin{table}[!htbp]
\caption{Функциональные слои фреймворка \code{atm}}
\label{tab:arch:layers}
\small
\begin{tabularx}{\textwidth}{l X}
\toprule
\textbf{Слой} & \textbf{Назначение} \\
\midrule
\code{core}          & Базовые типы и состояние графа --- сообщения, фазы,
                       роли, контейнеры состояния \code{AgentState},
                       \code{SharedState}, \code{GraphState}. \\
\code{llm}           & Унифицированная оболочка над провайдерами больших
                       языковых моделей, подсчет стоимости, бюджетные
                       ограничители, детерминированный \code{FakeLLM}
                       для тестов. \\
\code{tools}         & Реестр инструментов и песочница на~основе
                       контейнерной платформы Docker для
                       исполнения генерируемого кода. \\
\code{agents}        & Реализация пяти ролей агентов и базовый класс
                       \code{Agent} с управляющим циклом вызова
                       инструментов. \\
\code{topology}      & Шесть реализаций топологий поверх LangGraph
                       и общий протокол \code{Topology}. \\
\code{phases}        & Конечный автомат фаз, маршрутизатор топологий
                       и защитные правила переключения. \\
\code{human}         & Протокол \code{HumanGateway}, симулированный
                       и интерактивный шлюзы, маршрутизатор ролей
                       человека. \\
\code{storage}       & Объектно-реляционное отображение (от англ.
                       Object-Relational Mapping, ORM), асинхронные
                       сессии PostgreSQL, запись в формат Apache Parquet. \\
\code{observability} & Перехват событий LangGraph, сериализация
                       и отправка их в хранилище. \\
\code{tasks}         & Реализация и загрузчики бенчмарков
                       HumanEval, GSM8K, CommonGen, InfiAgent-DABench. \\
\code{evaluation}    & Судьи на основе языковой модели, оракулы
                       достоверности, агрегатор шкалы NASA-TLX
                       (от англ. NASA Task Load Index --- индекс
                       когнитивной нагрузки NASA). \\
\code{experiment}    & Запуск одиночных и грид-экспериментов,
                       схемы конфигов, командная строка. \\
\code{analysis}      & Загрузчики результатов, агрегаторы метрик,
                       построение графиков. \\
\bottomrule
\end{tabularx}
\end{table}

Архитектурный стиль опирается на три паттерна --- расширяемые точки
как протоколы Python (\code{Topology}, \code{HumanGateway},
\code{PhaseRouter}, \code{TopologyRouter}, \code{RoleRouter},
\code{Tool}) с \code{@runtime_checkable}; редукторное слияние
состояния графа при параллельном выполнении агентов с
идемпотентностью по идентификатору сообщения; реестры с декораторной
семантикой для топологий, инструментов и оракулов.

Технологический стек --- LangGraph~\cite{langgraph2024} поверх
LangChain~\cite{langchain2023} с~чекпойнтером для возобновления;
PostgreSQL~16 через асинхронные сессии SQL\-Alchemy~2.x для
метаданных и Apache Parquet для объемных журналов; YAML-конфиги
с~валидацией Pydantic~\cite{pydantic2024} поверх OmegaConf;
трассировка через LangSmith~\cite{langsmith2024}; линтер \code{ruff}
и анализатор типов \code{mypy} в строгом режиме.

Часть архитектурных решений заимствована из существующих работ
и фреймворков, часть является авторской разработкой. Точное
разграничение приведено в таблице~\ref{tab:arch:contribution}.

\begin{table}[!htbp]
\caption{Разделение архитектурных компонентов на заимствованные
и авторские}
\label{tab:arch:contribution}
\small
\begin{tabularx}{\textwidth}{>{\raggedright\arraybackslash}p{4.0cm}
                              >{\raggedright\arraybackslash}p{3.8cm}
                              >{\raggedright\arraybackslash}X}
\toprule
\textbf{Компонент} & \textbf{Источник} & \textbf{Авторский вклад} \\
\midrule
Оркестрация графов        & LangGraph (внешний) & Адаптация под мета-граф \\
Пять статичных топологий  & Прототипы из обзора литературы &
                                                     Унифицированная
                                                     реализация поверх
                                                     общего протокола \\
Адаптивная мета-топология & Своя разработка    & Полностью \\
Маршрутизаторы фаз        & Своя разработка    & Полностью, три режима
                                                 (rule-based, на основе
                                                 языковой модели, оракул) \\
Защитные правила переключения & Своя разработка & Полностью, четыре типа
                                                 ограничителей \\
Шлюз HITL                 & Своя разработка    & Полностью, три ролевых
                                                 маршрутизатора и политика
                                                 тайм-аутов \\
Оракул per-task best-of   & Своя разработка    & Полностью, в том числе
                                                 процедура построения
                                                 идеализированного референса \\
Подсистема наблюдаемости  & LangGraph callbacks & Адаптеры для PostgreSQL
                                                 и Apache Parquet, схема
                                                 событий \\
Бенчмарки                 & HumanEval, GSM8K, CommonGen, DABench &
                                                 Унифицированные загрузчики
                                                 и оценщики \\
Реализация шкалы NASA-TLX & Стандартная форма  & Прокси-агрегатор для
                                                 симулятора человека \\
\bottomrule
\end{tabularx}
\end{table}

\subsection{Слой агентов и оболочка над языковой моделью}
\label{sec:arch:agents}

Класс \code{Agent} в \code{atm.agents} реализует управляющий цикл ---
чтение входящих сообщений, формирование подсказки, вызов модели,
обработка инструментов и публикация ответа. Определены пять ролей
(\code{Planner}, \code{Researcher}, \code{Executor}, \code{Critic},
\code{Debater}) и дополнительный \code{Coordinator} для
иерархических и адаптивных топологий. Каждая роль задается
YAML-конфигом (подсказка, инструменты, окно сообщений, температура,
управление контекстом). Скрэтчпад --- скользящее окно с
автоматическим сжатием старых сообщений вспомогательной моделью;
политика выбрана по пилотному эксперименту.

Оболочка \code{LLMWrapper} в \code{atm.llm} унифицирует провайдеров
(OpenAI, Anthropic, Cerebras, локальный \code{FakeLLM}),
инкапсулирует подсчет токенов и стоимости (таблица \code{Pricing}),
экспоненциальный повтор и фиксацию версии модели. Ограничители
\code{BudgetGuard} работают на трех уровнях (вызов, запуск,
грид-эксперимент). \code{FakeLLM} поддерживает три режима ---
сценарный, эхо и воспроизведение по записи --- для детерминированного
повторения экспериментов.

\subsection{Реализация топологий поверх LangGraph}
\label{sec:arch:topology}

Каждая из шести топологий реализована классом, удовлетворяющим
протоколу \code{Topology} с фабричным методом \code{build},
возвращающим скомпилированный граф LangGraph. Граф связывает
узлы --- агентов и управляющие функции --- ребрами трех типов
(безусловные, условные, параллельные через примитив \code{Send}).
Все графы подключены к чекпойнтеру и перехватчику событий
\code{atm.observability}.

Пять статических топологий описаны в~\ref{sec:meth:topologies}.
Шестая, \emph{Adaptive}, --- мета-граф второго уровня. На каждом
тике маршрутизаторы фаз и топологий выбирают базовую топологию,
вызываемую как подграф; ворота \code{TransitionGate} применяют
правила передачи состояния и записывают \code{TopologyTransition}
в журнал. Подграфы кэшируются по имени для исключения повторной
компиляции; решения маршрутизаторов хранятся в замыкании, а не
в общем состоянии, чтобы избежать перезаписи во вложенном подграфе.

\subsection{Менеджер фаз и маршрутизатор топологий}
\label{sec:arch:phases}

Решение моделируется конечным автоматом (от англ. Finite-State
Machine, FSM) из четырех состояний (планирование, исполнение,
проверка, завершение). Маршрутизатор фаз в \code{atm.phases}
реализован в двух режимах --- \code{RuleBasedPhaseRouter} (таблица
сигналов с монотонной проверкой, фаза не движется назад) и
\code{LLMPhaseRouter} (парсинг ответа модели в формате
JSON --- от англ. JavaScript Object Notation --- с откатом на
правило при любой ошибке валидации).

Маршрутизатор топологий реализован в трех режимах. \code{RuleBased}
использует таблицу (фаза $\times$ сигналы)~$\to$~топология
(\code{stuck} в исполнении $\to$ Mesh, \code{rejected_count}~$\ge 3$
$\to$ Debate); пороги подобраны в пилотном эксперименте.
\code{LLMTopologyRouter} принимает решение по подсказке с описанием
состояния, с rule-based откатом. \code{OracleTopologyRouter} читает
таблицу решений из файла --- реализует идеализированный референс;
применяется только для построения идеализированного per-task best-of
референса и~в~экспериментальной воронке главы~\ref{sec:experiments}
как~самостоятельный режим маршрутизации не~запускается. Сравнение
RQ2 ведется между двумя реализуемыми в~условиях открытого развертывания
режимами~--- \code{RuleBased} и~\code{LLMTopologyRouter}.

Защитные правила (\code{SwitchGuards}) предотвращают
осцилляционный режим --- четыре ограничения: \code{min_dwell},
\code{cooldown}, лимиты за фазу и за запуск. Заблокированные
решения фиксируются как \code{guard_override} для последующей
оценки доли блокировок как метрики стабильности.

\subsection{Шлюз взаимодействия с человеком}
\label{sec:arch:human}

Интеграция человека абстрагирована протоколом \code{HumanGateway}
с асинхронным методом запроса ответа (идентификаторы запуска и
запроса, роль, история, допустимые действия). Идемпотентность по
паре (\code{run_id}, \code{request_id}) --- повторный вызов
возвращает ранее полученный ответ, что критично для возобновления
прерванных запусков.

Реализованы два шлюза --- \code{LLMSimulatedGateway} (языковая
модель с подсказкой по роли) и \code{CLIGateway} (интерактивный
интерфейс командной строки от англ. Command-Line Interface, CLI
для реальных участников). Маршрутизатор ролей выбирает
активную роль динамически (rule-based таблица или решение модели)
либо фиксированно (\code{FixedRoleRouter}). Политика тайм-аутов
имеет три варианта (откат к модели, пропуск, повтор).

\subsection{Хранение, командная строка и воспроизводимость}
\label{sec:arch:storage}

Метаданные --- в семи таблицах PostgreSQL (\code{experiments},
\code{runs}, \code{phases}, \code{topology_transitions},
\code{human_interactions}, \code{budget_events},
\code{llm_calls}) с индексами по типичным аналитическим запросам
и составным индексом (эксперимент, статус) для выборки
невыполненных запусков. Объемные журналы --- в файлах Apache
Parquet с асинхронным писателем и партиционированием по запуску.
Миграции --- Alembic.

CLI \code{atm} на библиотеке Typer объединяет семь команд
(таблица~\ref{tab:arch:cli}).

\begin{table}[!htbp]
\caption{Команды интерфейса командной строки фреймворка}
\label{tab:arch:cli}
\small
\begin{tabularx}{\textwidth}{l X}
\toprule
\textbf{Команда} & \textbf{Назначение} \\
\midrule
\code{atm run}       & Запуск одиночного эксперимента из YAML-конфига. \\
\code{atm grid}      & Параллельный грид-эксперимент через
                       \code{ProcessPoolExecutor}. \\
\code{atm estimate}  & Предварительная оценка стоимости запуска
                       по таблице цен и историческим данным. \\
\code{atm status}    & Агрегированный статус эксперимента ---
                       число выполненных, неудачных и активных
                       ячеек, суммарная стоимость. \\
\code{atm resume}    & Возобновление прерванного запуска из
                       последнего чекпойнта LangGraph. \\
\code{atm replay}    & Детерминированное воспроизведение запуска
                       через \code{FakeLLM} в режиме воспроизведения
                       по записи. \\
\code{atm reconcile} & Поиск и пометка запусков-зомби (запись о статусе
                       \enquote{выполняется}, но процесс уже
                       завершен). \\
\bottomrule
\end{tabularx}
\end{table}

Воспроизводимость --- инварианты, фиксируемые при старте запуска
(\code{seed_all} для Python/NumPy/провайдеров/генератора
идентификаторов; \code{git_sha} из \code{git rev-parse HEAD};
снимок версий моделей в \code{model_version_snapshot} по первому
ответу провайдера; дайджест образа Docker-песочницы). Совокупность
полей делает возможным сравнение запусков между собой и анализ
артефактов депрекации моделей.

\subsection{Выводы и результаты по главе}
\label{sec:arch:summary}

Инфраструктура реализована как фреймворк из тринадцати слоев на
Python с расширяемостью через протоколы, воспроизводимостью через
фиксацию инвариантов и наблюдаемостью через единый перехват
событий LangGraph; в нее входят шесть топологий, два режима фаз,
три режима топологий и два шлюза HITL --- все необходимое для
экспериментов главы~\ref{sec:experiments}. Разграничение
заимствованных и авторских компонентов --- в
таблице~\ref{tab:arch:contribution}, авторский вклад
развернут в~\ref{sec:contribution}. Кодовая база основного пакета
\code{atm} --- 21\,327 строк Python в 112 модулях внутри
тринадцати слоев; конфигурация --- 43 YAML-файла (24 описывают
запуски E1--E4).
